%% kapa: exam=1 round=31 year=2020 period=1 form=A subject=real_estate subject_ko=부동산학원론 seq=19 number=99
%% source: Q-Net 감정평가사 기출문제 (cst00309) artlSeq=5207219 file=extracted/1008/5207219/A형/A형/2020년도 제31회 감정평가사 1차 시험_1교시_A형.pdf
%% conversion: kordoc 4.24.51 -> tools/gen1.py (auto); review=proofread-sonnet; flags=none
\begin{question}{99}
건물의 관리방식에 관한 설명으로 옳지 않은 것은?
\begin{choices}
  \item 자가관리방식은 일반적으로 소유자의 지시와 통제 권한이 강하다.
  \item 위탁관리방식은 부동산관리를 전문적으로 하는 대행업체에게 맡기는 방식으로 사회적으로 신뢰도가 높고 성실한 대행업체를 선정하는 것이 중요하다.
  \item 혼합관리방식은 자가관리에서 위탁관리로 이행하는 과도기적 조치로 적합하다.
  \item 자가관리방식에 있어 소유자가 전문적 관리지식이 부족한 경우 효율적 관리에 한계가 있을 수 있다.
  \item 혼합관리방식에 있어 관리상의 문제가 발생할 경우, 책임소재에 대한 구분이 명확하다.
\end{choices}
\end{question}
\answer{5}
